\documentclass[10pt]{article}
\usepackage[margin=0.8in]{geometry}
\usepackage{amsmath,amsfonts}
\providecommand{\IEEEQED}{\rule{1.3ex}{1.3ex}}
\setlength{\parskip}{3pt}
% Paper equation numbers, copied from the revision's .aux so references resolve here.
\makeatletter
\newlabel{eq:2}{{1}{3}}\newlabel{eq:3}{{2}{3}}\newlabel{eq:9}{{6}{3}}\newlabel{eq:scope}{{7}{3}}
\newlabel{eq:10}{{8}{3}}\newlabel{eq:11}{{9}{3}}\newlabel{eq:momentum}{{11}{4}}
\makeatother
\begin{document}
\setcounter{equation}{11}

\noindent\textbf{Proposed Section III-D, placed after the robot model so that (\ref{eq:momentum}) is defined before use.}

\subsection*{D. Closed-Loop Stability and Robustness}

The convergence results of Section~II-D hold for a stationary estimate and an ideal centroid. In operation the estimate changes every cycle with the robot positions and sensor noise, and the centroid responds through the sampled lag. At cycle $j$, write the estimate as $\hat{\mathbf{p}}^*[j] = \mathbf{p}^* + \boldsymbol{\delta}[j]$ and the centroid error as $\mathbf{e}[j] = \mathbf{p}_c[j] - \mathbf{p}^*$. If the formation is held, the centroid follows (\ref{eq:momentum}) and advances as $\mathbf{p}_c[j+1] = \mathbf{p}_c[j] + \Delta t\,\mathbf{v}[j+1]$.

In continuous time, (\ref{eq:10}) gives $\dot{\mathbf{e}} = -k(\mathbf{e} - \boldsymbol{\delta})$. Sampled with lag, each axis of $\mathbf{x} = [e,\ v]^T$ evolves as
\begin{equation}\label{eq:cl}
\mathbf{x}[j+1] = \begin{bmatrix} 1-(1-\alpha)k\Delta t & \alpha\Delta t \\ -(1-\alpha)k & \alpha \end{bmatrix}\mathbf{x}[j] + (1-\alpha)k\begin{bmatrix}\Delta t\\ 1\end{bmatrix}\delta[j].
\end{equation}
The characteristic polynomial $z^2 - (1+\alpha-(1-\alpha)k\Delta t)\,z + \alpha$ meets the Jury conditions, and the loop is exponentially stable, if and only if
\begin{equation}\label{eq:kmax}
0 < k\Delta t < 2(1+\alpha)/(1-\alpha).
\end{equation}
The impulse response from $\delta$ to $e$ sums to one. Its $\ell_1$ norm $\gamma$ equals one when both poles are real and positive ($k\Delta t \le (1-\sqrt{\alpha})/(1+\sqrt{\alpha})$), and at our operating point ($\Delta t = 0.1$~s, $\alpha = 0.70$, $k = 1$~s$^{-1}$) it is $\gamma = 1.0003$.

Coulomb friction removes at most $(1-\alpha)v_{\min}$ of speed per cycle and holds a parked robot while its command is below $v_{\min}$. Both effects equal a command disturbance of norm at most $v_{\min}$, which (\ref{eq:cl}) receives as an extra estimate error of at most $v_{\min}/k$. Per-axis saturation bounds every velocity by $v_{\max}$ whatever the estimate. The estimate error itself depends on cluster position. By (\ref{eq:scope}), $\|\boldsymbol{\delta}\| \le \|\hat{\mathbf{v}}(\mathbf{p}^*) - \mathbf{v}(\mathbf{p}^*)\|/\sigma_{\min}(\hat{\mathbf{J}})$, and by (\ref{eq:2})--(\ref{eq:3}) the numerator grows with field curvature as $(\|\mathbf{e}\| + R)^2$, where $R$ is the formation circumradius, plus sensor noise extrapolated from the formation to $\mathbf{p}^*$.

\smallskip
\noindent\textbf{Proposition 2 (Bounded tracking):} \textit{Assume the formation is held, (\ref{eq:kmax}) holds, the command near rest is unsaturated, and on a bounded region containing the trajectory $\|\boldsymbol{\delta}[j]\| \le \eta\|\mathbf{e}[j]\| + \bar{\delta}$ with $\gamma\eta < 1$. Then}
\begin{equation}\label{eq:ultimate}
\limsup_{j\to\infty} \|\mathbf{e}[j]\| \le \frac{\gamma\,(\bar{\delta} + v_{\min}/k)}{1 - \gamma\eta}.
\end{equation}

\noindent\textit{Proof:} Under (\ref{eq:kmax}), (\ref{eq:cl}) is exponentially stable, so its output obeys $\limsup\|\mathbf{e}\| \le \gamma\limsup\|\boldsymbol{\delta}_{\mathrm{eff}}\|$ for the effective input $\|\boldsymbol{\delta}_{\mathrm{eff}}\| \le \|\boldsymbol{\delta}\| + v_{\min}/k$. Substituting the bound on $\boldsymbol{\delta}$ gives $L \le \gamma(\eta L + \bar{\delta} + v_{\min}/k)$ for the finite $L = \limsup\|\mathbf{e}\|$, which rearranges to (\ref{eq:ultimate}) when $\gamma\eta < 1$. \hfill\IEEEQED

\smallskip
Curvature, noise, and friction therefore enlarge the final error without destabilizing the loop. Both $\eta$ and $\bar{\delta}$ scale with $1/\sigma_{\min}(\hat{\mathbf{J}})$, so Proposition~2 gives no guarantee when $\hat{\mathbf{J}}$ is nearly singular, as in near-uniform flow far from a critical point. Saturation still caps the centroid speed at $v_{\max}$ per axis, and the memoryless estimator replaces a spurious estimate on the next cycle. The experiments reported here applied no conditioning gate to (\ref{eq:9}).

For the orbital law, let $r^* = \|\mathbf{p}^* - \mathbf{p}_c\|$ be the true radius and split $\boldsymbol{\delta}$ into components $\delta_n$ along $\mathbf{n}$ and $\delta_t$ along $\mathbf{t}$. To first order in $\|\boldsymbol{\delta}\|/r^*$, $r^*$ obeys
\begin{equation}\label{eq:orbit_iss}
\dot{r}^* = -k_r(r^* - r_d) - k_r\delta_n + (k_t/r^*)\,\delta_t,
\end{equation}
so the radial error is bounded by $\sup|\delta_n| + k_t\sup|\delta_t|/(k_r r_d)$ in addition to the steady outward offset of lag and sampling. A tangential estimate error is amplified by $k_t/(k_r r_d)$, which grows without limit at small radii. The singular point $r = 0$ of (\ref{eq:11}) is not reached by the ideal dynamics, since $r(t) = r_d + (r(0) - r_d)e^{-k_r t} > 0$. A noisy estimate can still land on $\mathbf{p}_c$, and any fixed choice of $\mathbf{n}$ there, such as $[1, 0]^T$, lets the tangential step $k_t\mathbf{t}$ carry the cluster away on the next cycle.

\bigskip
\noindent\textbf{Section VI-A, last paragraph, replaced by:} Across all four error sources, the estimation error enters the closed loop as the input $\boldsymbol{\delta}$ of (\ref{eq:cl}). Sensor noise, curvature, and conditioning enlarge the final attraction error by (\ref{eq:ultimate}) and the radial orbit error by (\ref{eq:orbit_iss}) without destabilizing either controller, and the guarantee lapses only as $\hat{\mathbf{J}}$ approaches singularity.

\noindent\textbf{Section VI-A, same section:} ``non-hyperbolic critical points'' becomes ``degenerate critical points''.

\end{document}
